\section{实验原理}

\subsection{针孔双目模型}

设左右相机光心在水平方向相距基线 $B$，两相机具有相同焦距 $f$。对空间点
$(X,Y,Z)$，在平行相机模型下，其水平投影可写为
\begin{equation}
  x_L=\frac{fX}{Z},\qquad
  x_R=\frac{f(X-B)}{Z}.
\end{equation}
两式相减得到视差与深度关系
\begin{equation}
  d=x_L-x_R=\frac{fB}{Z},\qquad
  Z=\frac{fB}{d}.
  \label{eq:metric-depth}
\end{equation}
因此，较大的视差通常对应较近的场景点；较小的视差对应较远的场景点。式~\eqref{eq:metric-depth}
还说明，像素视差本身不能决定深度的物理尺度，必须同时知道与当前相机系统匹配的
焦距和基线。

\subsection{对极约束与水平搜索}

双目图像中的对应点满足基础矩阵给出的对极约束
\begin{equation}
  \mathbf{x}_R^{\mathsf{T}}\mathbf{F}\mathbf{x}_L=0,
\end{equation}
其中 $\mathbf{x}_L$、$\mathbf{x}_R$ 为对应齐次像素坐标，$\mathbf{F}$ 为基础矩阵。
该关系把右图中的候选点限制在左图点对应的对极线上。经过可靠的立体校正后，对极线
近似变为水平扫描线，因此匹配器只需在同一行内搜索不同的水平位移。

本实验沿用 OpenCV aloe 官方样例的直接匹配流程，不重新估计校正变换。报告中的视差
约定为 $d=x_L-x_R$；程序输出的有效视差规则是有限值且 $d>15$。

\subsection{StereoBM 与 StereoSGBM}

\textbf{StereoBM} 在局部窗口内比较左右图像块，结构简单、速度较快，但在纹理不足、
遮挡边界和重复纹理处容易出现空洞或错误匹配。\textbf{StereoSGBM} 在多个方向上累积
匹配代价，使用半全局的路径优化在局部一致性和全局平滑之间折中，通常能得到更连续
的视差结果。本实验选择 StereoSGBM，以便在官方 aloe 图像上获得可解释的课程演示结果。

对像素 $p$ 和候选视差 $d$，可将匹配代价抽象为 $C(p,d)$。沿某条路径的递推会惩罚
相邻像素的视差变化：视差变化较小时使用惩罚 $P_1$，发生更大跳变时使用惩罚 $P_2$。
本实验使用 $P_2>P_1$，鼓励大多数区域保持平滑，同时允许物体边界出现视差突变。实际
配置还包含以下约束：

\begin{itemize}
  \item \texttt{minDisparity} 与 \texttt{numDisparities} 定义搜索起点和搜索范围；本次为 $16$ 与 $96$；
  \item \texttt{blockSize} 决定局部匹配窗口，本次为 $5$；窗口越大通常越稳定，但会损失边界细节；
  \item \texttt{uniquenessRatio} 要求最佳匹配明显优于次佳匹配，降低歧义点被保留的概率；
  \item \texttt{disp12MaxDiff} 用于左右一致性检查；
  \item \texttt{speckleWindowSize} 与 \texttt{speckleRange} 去除小面积、视差变化异常的散斑区域。
\end{itemize}

\subsection{相对深度与标定边界}

OpenCV 匹配器返回带固定缩放的整数结果，程序除以 $16$ 得到像素单位的浮点视差。
对有效视差计算
\begin{equation}
  \mathrm{depth\_proxy}=\frac{1}{d},\qquad d>15.
  \label{eq:proxy-depth}
\end{equation}
该量保留视差反映的前后顺序，但不包含式~\eqref{eq:metric-depth} 中的 $fB$。
当前官方 aloe 样例没有明确配套的 $f$、$B$、\texttt{doffs} 或 $\mathbf{Q}$，因此不调用带猜测
参数的 \texttt{reprojectImageTo3D}，也不生成带物理尺度的三维点云。\texttt{aloeGT.png} 如被使用，
只能作为视差参考图，不能作为米制深度真值。

\subsection{视差误差的深度敏感性}

式~\eqref{eq:metric-depth} 不仅给出了深度的反演公式，还给出了误差传播方向。对
固定的 $fB$ 求导可得
\begin{equation}
  \frac{\partial Z}{\partial d}=-\frac{fB}{d^2},\qquad
  \frac{\Delta Z}{Z}\approx-\frac{\Delta d}{d}.
  \label{eq:depth-sensitivity}
\end{equation}
对于本报告使用的相对深度代理，也有
\begin{equation}
  \frac{\partial (1/d)}{\partial d}=-\frac{1}{d^2}.
\end{equation}
这意味着相同大小的像素视差扰动，在小视差区域会造成更大的相对影响；而在大视差区域，
虽然场景通常更近，但单个像素误差对 $1/d$ 的绝对变化较小。因此，不能只依据视差图的
灰度连续性判断每个像素的深度精度，还应结合有效掩码、局部纹理和匹配一致性共同解释。

\subsection{有效视差的定义与统计含义}

程序把有限且满足 $d>15$ 的像素定义为有效视差，其他位置保留为 \texttt{NaN}。这里的
阈值首先是数值和搜索范围约束：匹配器以 $16$ 为搜索起点，低于阈值的值不能直接当作
可信的正视差。该规则回答的是“哪些像素进入后续计算”，并不等价于“这些像素全部匹配
正确”。因此，有效掩码是后续统计的支持集，而不是精度标签；有效像素比例是覆盖率，
不是准确率。

在缺少可靠视差真值时，至少应把三类信息分开记录：一是覆盖率，即有多少像素通过有效性
筛选；二是视差分布，即有效像素落在什么范围和分位数区间；三是匹配可信度，即这些视差
是否通过左右一致性、重投影误差或外部参考的检验。本实验保存原始 \texttt{npy} 数组和
掩码，正是为了避免将经过百分位拉伸的显示 PNG 误当作上述定量证据。

\subsection{下采样后的坐标与可比性}

本次运行对左右图像同步执行一次 \texttt{pyrDown}，匹配坐标位于 $641\times555$ 的处理
图像上，而不是原始 $1282\times1110$ 图像上。因而，视差数组中的一个像素对应处理图像
坐标中的一个像素；若将结果与原始分辨率下的标定参数或另一组未下采样结果比较，必须
同步考虑焦距、主点和视差坐标的尺度变换。否则，即使场景和算法完全相同，也可能把
分辨率变化造成的数值差异误读为深度变化。

下采样的直接收益是减少匹配计算量，并使课程实验可以快速复现；代价是细小纹理和窄边界
可能被平滑或合并。由于本实验只在同一处理分辨率内解释 $d$ 和 $1/d$ 的相对关系，报告
不把该结果外推为原始相机坐标系下的米制深度。若未来启用标定参数，应在同一分辨率下
重新生成标定模型，或明确地将参数换算到处理分辨率。
